\section{Environnement de développement}

\subsection{Outils et choix technologiques}

Le développement est effectué sur un poste Windows avec PHP installé localement
et Docker Desktop appuyé sur WSL2 pour les composants conteneurisés. Cette
organisation conserve un cycle rapide pour l'interface et l'API tout en isolant
la passerelle OCPP, le simulateur de bornes et le simulateur de paiement. Git
assure la traçabilité du code et GitHub centralise le dépôt distant.

\begin{table}[H]
  \centering
  \caption{Environnement technique effectivement utilisé}
  \label{tab:implementation-toolchain}
  \begin{tabularx}{\textwidth}{|p{3.2cm}|p{4.0cm}|Y|}
    \hline
    \rowcolor{ctmint}
    \textbf{Domaine} & \textbf{Technologies} & \textbf{Usage dans le projet}\\\hline
    Frontend & React 19, TypeScript 6, Vite 8 &
    Application monopage, typage statique et compilation.\\\hline
    Interface & Ant Design 6, Recharts, Leaflet, GSAP &
    Composants, graphiques, cartes et animations ciblées.\\\hline
    État et échanges & TanStack Query, Axios, Zustand &
    Cache serveur, appels HTTP et états locaux transverses.\\\hline
    Backend & PHP 8.5, Laravel 13 &
    API REST, logique métier, validation, autorisation et tâches planifiées.\\\hline
    Identité & Sanctum, Socialite, Spatie Permission &
    Session SPA, Google OAuth 2.0, rôles et permissions.\\\hline
    Données & PostgreSQL 16, Eloquent &
    Persistance transactionnelle, relations et migrations.\\\hline
    Temps réel & Reverb, Laravel Echo &
    Notifications et mises à jour poussées par WebSocket.\\\hline
    OCPP & Python 3.11, OCPP 1.6J, simulateur SAP &
    Passerelle WebSocket, messages de borne et scénarios de test.\\\hline
    Paiement & Contrat \texttt{PaymentGateway}, WireMock 3 &
    Fournisseur externe simulé, webhooks et scénarios d'erreur.\\\hline
    Email & Resend, queue Laravel &
    Vérification d'email, invitations et notifications asynchrones.\\\hline
    Qualité & PHPUnit 12, Oxlint, TypeScript &
    Tests backend, analyse frontend et validation de compilation.\\\hline
    Documentation & \LaTeX, TikZ, OpenAPI 3.1 avec Scramble &
    Rapport, schémas et spécification versionnée de l'API.\\\hline
  \end{tabularx}
\end{table}

Les dépendances sont verrouillées par \texttt{composer.lock} et
\texttt{pnpm-lock.yaml}. Les valeurs sensibles sont placées dans les fichiers
\texttt{.env}, exclus de Git. Les fichiers \texttt{.env.example} documentent
uniquement les variables attendues, sans contenir de secret.

\subsection{Langages de programmation utilisés}

\begin{table}[H]
  \centering
  \caption{Langages utilisés dans la réalisation}
  \label{tab:implementation-languages}
  \begin{tabularx}{\textwidth}{|p{3.2cm}|p{3.8cm}|Y|}
    \hline
    \rowcolor{ctmint}
    \textbf{Langage} & \textbf{Composant} & \textbf{Responsabilité}\\\hline
    PHP & Backend Laravel &
    API REST, services métier, jobs, policies et tests.\\\hline
    TypeScript / TSX & Frontend React &
    Pages, composants, contrats API et interactions utilisateur.\\\hline
    Python & Passerelle OCPP &
    WebSocket OCPP, adaptation des messages et tests du protocole.\\\hline
    SQL & PostgreSQL &
    Migrations, contraintes, index et requêtes d'agrégation.\\\hline
  \end{tabularx}
\end{table}

\subsection{Organisation du dépôt et environnement d'exécution}

Le dépôt sépare physiquement les composants tout en proposant des commandes
communes à sa racine :

\begin{description}[leftmargin=37mm,style=nextline]
  \item[\texttt{frontend/}] application React et ressources visuelles ;
  \item[\texttt{backend/}] API Laravel, services, migrations, policies et tests ;
  \item[\texttt{ocpp-gateway/}] passerelle Python entre les bornes et Laravel ;
  \item[\texttt{infra/ocpp/}] composition Docker et profils du simulateur SAP ;
  \item[\texttt{infra/payment-simulator/}] serveur WireMock et contrats simulés ;
  \item[\texttt{scripts/}] commandes d'orchestration et de pilotage ;
  \item[\texttt{docs/}] rapport, API, manuels et cahier de tests.
\end{description}

\begin{figure}[H]
  \centering
  \resizebox{0.94\textwidth}{!}{\begin{tikzpicture}[
  font=\sffamily\small,
  process/.style={
    draw=ctline, rounded corners=2mm, thick, fill=white,
    minimum width=34mm, minimum height=13mm, align=center
  },
  hostproc/.style={process, draw=ctgreen, fill=ctmint},
  container/.style={process, draw=ctblue, fill=blue!5},
  optional/.style={process, dashed, draw=ctgray, fill=gray!4},
  flow/.style={{Latex[length=2.2mm]}-{Latex[length=2.2mm]}, thick, draw=ctgray},
  zone/.style={draw=ctline, rounded corners=3mm, dashed, inner sep=6mm}
]
  \node[hostproc] (vite) at (0,2.2) {Vite\\\texttt{:5173}};
  \node[hostproc] (laravel) at (5.0,2.2) {Laravel API\\\texttt{:8000}};
  \node[hostproc] (reverb) at (10.0,2.2) {Reverb\\\texttt{:8080}};
  \node[hostproc] (worker) at (2.5,-0.1) {Queue worker};
  \node[hostproc] (scheduler) at (7.5,-0.1) {Scheduler};

  \node[container] (postgres) at (0,-4.2) {PostgreSQL\\\texttt{:5432}};
  \node[optional] (mailpit) at (5.0,-4.2) {Mailpit optionnel\\\texttt{:8025/1025}};
  \node[container] (wiremock) at (10.0,-4.2) {WireMock\\\texttt{:9090}};
  \node[container] (gateway) at (0,-6.8) {Gateway OCPP\\\texttt{:9000}};
  \node[container] (sap) at (5.0,-6.8) {Simulateur SAP\\\texttt{:8082}};

  \draw[flow] (vite) -- (laravel);
  \draw[flow] (vite) to[bend left=18] (reverb);
  \draw[flow] (laravel) -- (reverb);
  \draw[flow] (laravel) -- (worker);
  \draw[flow] (laravel) -- (scheduler);
  \draw[flow] (laravel.south) to[out=-115,in=75] (postgres.north);
  \draw[dashed, flow] (worker.south) to[out=-35,in=160] (mailpit.west);
  \draw[flow] (laravel.south) to[out=-145,in=70] (gateway.north);
  \draw[flow] (gateway) -- (sap);
  \draw[flow] (laravel.south east) to[out=-35,in=95] (wiremock.north);

  \node[zone, fit=(vite)(laravel)(reverb)(worker)(scheduler),
    label={[ctgreen]above:\textbf{Processus hôte}}] {};
  \node[zone, fit=(postgres)(mailpit)(gateway)(sap)(wiremock),
    label={[ctblue]above:\textbf{Docker Desktop / WSL2}}] {};
\end{tikzpicture}
}
  \caption{Architecture de déploiement de l'environnement local}
  \label{fig:deployment-architecture}
\end{figure}

En production, les fichiers frontend doivent être servis derrière HTTPS,
l'application PHP doit être supervisée et les connexions OCPP doivent utiliser
\texttt{wss://}. Les workers, le scheduler, Reverb et la passerelle OCPP
constituent des processus distincts à superviser.

\subsection{Processus de développement}

Le fichier \texttt{package.json} racine sert de façade aux commandes des
différents environnements :

\begin{table}[H]
  \centering
  \caption{Commandes principales de l'environnement local}
  \label{tab:implementation-commands}
  \begin{tabularx}{\textwidth}{|p{4.4cm}|Y|}
    \hline
    \rowcolor{ctmint}
    \textbf{Commande} & \textbf{Responsabilité}\\\hline
    \texttt{pnpm dev:frontend} & Lance Vite sur le port 5173.\\\hline
    \texttt{pnpm dev:backend} & Lance l'API Laravel sur le port 8000.\\\hline
    \texttt{pnpm dev:queue} & Traite les emails et travaux asynchrones.\\\hline
    \texttt{pnpm dev:scheduler} & Exécute les règles périodiques.\\\hline
    \texttt{pnpm dev:reverb} & Lance le serveur WebSocket applicatif.\\\hline
    \texttt{pnpm ocpp:up} & Démarre la passerelle et le simulateur OCPP.\\\hline
    \texttt{pnpm ocpp:plug} / \texttt{ocpp:unplug} &
    Simule la connexion ou le retrait physique d'un câble.\\\hline
    \texttt{pnpm payment:up} & Démarre le fournisseur WireMock.\\\hline
    \texttt{pnpm build:frontend} & Vérifie les types et compile le frontend.\\\hline
    \texttt{pnpm test:backend} & Exécute la campagne PHPUnit.\\\hline
  \end{tabularx}
\end{table}

Le scheduler recalcule périodiquement la disponibilité, expire les commandes
OCPP, contrôle les SLA et génère les maintenances récurrentes. L'option
\texttt{withoutOverlapping} empêche le chevauchement d'une même tâche.

Chaque bloc fonctionnel est suivi d'un commit cohérent puis poussé sur la branche
\texttt{main}. Cette discipline facilite l'identification d'une régression. Les
fichiers d'environnement, les dépendances installées et les sorties temporaires
ne sont pas versionnés.
